\section{Theoretical Setup}

\subsection{Smoothed CP target and objective}

Let $r,n,k$ be positive integers and let $q>0$ be fixed.  For arbitrary
deterministic matrices
$\bar A=[\bar a_1,\ldots,\bar a_r]$,
$\bar B=[\bar b_1,\ldots,\bar b_r]$, and
$\bar C=[\bar c_1,\ldots,\bar c_r]$ in $\mathbb R^{n\times r}$, set
$\rho=r^{-q}$.  Independently draw
\[
 \xi_j^a,\xi_j^b,\xi_j^c\sim
 \mathcal N\!\left(0,\frac{\rho^2}{n}I_n\right),
 \qquad 1\le j\le r,
\]
and define
\[
 a_j=\bar a_j+\xi_j^a,\qquad
 b_j=\bar b_j+\xi_j^b,\qquad
 c_j=\bar c_j+\xi_j^c,
 \qquad
 T=\sum_{j=1}^r a_j\otimes b_j\otimes c_j.
\]
No conditioning, incoherence, or random-base restriction is imposed on the
deterministic matrices.

For $X=[x_1,\ldots,x_k]$, $Y=[y_1,\ldots,y_k]$, and
$Z=[z_1,\ldots,z_k]$ in $\mathbb R^{n\times k}$, write
\[
 S(X,Y,Z)=\sum_{i=1}^k x_i\otimes y_i\otimes z_i,
 \qquad
 F(X,Y,Z)=\frac12\|T-S(X,Y,Z)\|_F^2.
\]
The symbol $T_{(m)}$ denotes the mode-$m$ matricization, and $\odot$
denotes the Khatri--Rao product in one fixed standard ordering.

\subsection{Method-specific fixed initialization spans}

For each $M\in\{\mathrm{cALS},\mathrm{cGD}\}$, independently draw
$G_x^M,G_y^M,G_z^M\in\mathbb R^{n\times k}$ with iid
$\mathcal N(0,1/n)$ entries.  Let
\[
 \mathcal S_M=\operatorname{range}(G_x^M),
 \qquad
 Q_M=\operatorname{orth}(G_x^M),
\]
where $\operatorname{orth}$ is a fixed measurable choice of an orthonormal
basis on the full-column-rank set.  Define
\[
 \mathcal H_M=\mathcal S_M\otimes\mathbb R^n\otimes\mathbb R^n,
 \qquad
 P_{\mathcal S_M}=Q_MQ_M^{\mathsf T},
 \qquad
 P_{\mathcal H_M}=P_{\mathcal S_M}\otimes I_n\otimes I_n.
\]
The two initialization triples are independent of one another and of the
smoothing variables.  Both methods use the same realization of $T$.

The constrained sequential alternating least-squares method initializes
\[
 (X_0,Y_0,Z_0)=
 (G_x^{\mathrm{cALS}},G_y^{\mathrm{cALS}},G_z^{\mathrm{cALS}}).
\]
At sweep $t$, it updates in the order $X,Y,Z$.
With
\[
 K_t^x=Z_t\odot Y_t,
\]
the constrained block is
\[
 X_{t+1}=Q_{\mathrm{cALS}}
 \left[Q_{\mathrm{cALS}}^{\mathsf T}T_{(1)}K_t^x
 \bigl((K_t^x)^{\mathsf T}K_t^x\bigr)^\dagger\right].
\]
It is followed by
\[
 K_t^y=Z_t\odot X_{t+1},
 \qquad
 Y_{t+1}=T_{(2)}K_t^y
 \bigl((K_t^y)^{\mathsf T}K_t^y\bigr)^\dagger,
\]
and
\[
 K_t^z=Y_{t+1}\odot X_{t+1},
 \qquad
 Z_{t+1}=T_{(3)}K_t^z
 \bigl((K_t^z)^{\mathsf T}K_t^z\bigr)^\dagger.
\]
The displayed choices are the minimum-Frobenius-norm exact block
least-squares solutions, including at singular Gram matrices.  The method has
no restart or early stopping.

The constrained gradient method writes $X_t=Q_{\mathrm{cGD}}C_t$ and uses
\[
 f_{Q_{\mathrm{cGD}}}(C,Y,Z)=F(Q_{\mathrm{cGD}}C,Y,Z).
\]
It initializes
\[
 C_0=Q_{\mathrm{cGD}}^{\mathsf T}G_x^{\mathrm{cGD}},
 \qquad Y_0=G_y^{\mathrm{cGD}},
 \qquad Z_0=G_z^{\mathrm{cGD}}.
\]
At a current state $u_t=(C_t,Y_t,Z_t)$, it tests the dyadic step sizes
$\eta=2^{-j}$, $j=0,1,2,\ldots$, in order, and chooses the first one for
which, with $V_t=u_t-\eta\nabla f_{Q_{\mathrm{cGD}}}(u_t)$,
\[
 f_{Q_{\mathrm{cGD}}}(V_t)
 \le f_{Q_{\mathrm{cGD}}}(u_t)
 -\frac{\eta}{2}\|\nabla f_{Q_{\mathrm{cGD}}}(u_t)\|_F^2.
\]
Here the squared gradient norm is the sum of the squared Frobenius norms of
the $C,Y,Z$ blocks.  The accepted value is denoted $\eta_t$, and
$u_{t+1}=V_t$, $X_{t+1}=Q_{\mathrm{cGD}}C_{t+1}$.  This method also has no
restart.

For either method, set
\[
 S_t^M=S(X_t^M,Y_t^M,Z_t^M),
 \qquad
 F_M(t)=F(X_t^M,Y_t^M,Z_t^M).
\]
When $\|(I-P_{\mathcal H_M})T\|_F>0$, define the fixed normalized tensor
\[
 W_M=\frac{(I-P_{\mathcal H_M})T}
 {\|(I-P_{\mathcal H_M})T\|_F}.
\]
No value is assigned to $W_M$ when the denominator is zero.

\subsection{Assumptions}

\begin{assumption}[Fixed ambient dimension]\label{assump:dimension}
$r$ and $n$ are positive integers and $n\ge 8r^{5/4}$.
\end{assumption}

\begin{assumption}[Superlinear algorithmic rank]\label{assump:rank_window}
$k$ is an integer satisfying
\[
 r<k\le r^{1+c}=r^{5/4},
 \qquad c=\frac14.
\]
\end{assumption}

\begin{assumption}[Uniform arbitrary deterministic bases]\label{assump:arbitrary_base}
$\bar A,\bar B,\bar C\in\mathbb R^{n\times r}$ are deterministic and
unrestricted.  In particular, no conditioning or incoherence hypothesis is
imposed, and the theorem is pointwise uniform over all such triples.
\end{assumption}

\begin{assumption}[Independent Gaussian smoothing]\label{assump:gaussian_smoothing}
$q>0$ is fixed, $\rho=r^{-q}$, and all $3r$ perturbation vectors used above
are mutually independent with covariance $(\rho^2/n)I_n$.
\end{assumption}

\begin{assumption}[Shared target and independent starts]\label{assump:joint_initialization}
The two method-specific Gaussian initialization triples are mutually
independent and independent of the smoothing variables.  Both methods use the
same realized tensor $T$ when their events are intersected.
\end{assumption}

\subsection{Formalized goal}

The objective is to prove that one may take $r_0=1$ and that, for every
admissible $r,n,k$ and every unrestricted deterministic base triple, the joint
smoothing-and-initialization law satisfies
\[
 \mathbb P\!\left[
 \bigcap_{M\in\{\mathrm{cALS},\mathrm{cGD}\}}
 \left\{
 \lim_{t\to\infty}F_M(t)\text{ exists and }
 \lim_{t\to\infty}F_M(t)\ge\frac38\|T\|_F^2
 \right\}
 \right]\ge\frac14.
\]
This is a material-partial result for the one-mode fixed-span constrained
class.  Replacing the fixed initialization span by the endogenous moving span
of ordinary unconstrained ALS or full-variable gradient descent remains open.
